\par\Needspace{6\baselineskip}\medskip
\begin{center}\hypertarget{native-199-article149}{}\pdfbookmark[2]{149.}{article-149}149.\end{center}
Erst in der letzten Hypothese berechnet man die Bahnelemente selbst, und zwar entweder aus $f,r',r''$ oder aus $f'',r,r'$. Dazu führt man eine der beiden Rechnungen zu Ende, die man in den vorangegangenen Hypothesen nur bis zu $\eta$ oder $\eta''$ verfolgen musste. Führt man beide Rechnungen aus, so liefert die Übereinstimmung der erhaltenen Zahlen eine weitere Kontrolle der gesamten Arbeit. Sobald $f,f',f''$ gefunden sind, leitet man die Elemente jedoch zweckmäßiger allein aus der Kombination des ersten und dritten Ortes ab, also aus $f',r,r''$ und dem Zeitintervall. Zur besseren Kontrolle der Rechnung bestimmt man schließlich den mittleren Ort in der Bahn aus den gefundenen Elementen.

So erhält man die Abmessungen des Kegelschnitts, nämlich die Exzentrizität, die große Halbachse oder den Halbparameter und die Lage des Perihels zu den heliozentrischen Orten $C,C',C''$, ferner bei einer elliptischen Bahn die mittlere Bewegung und die mittlere Anomalie für eine beliebige Epoche oder bei einer hyperbolischen oder parabolischen Bahn den Zeitpunkt des Periheldurchgangs. Es bleibt also nur noch, die Lage der heliozentrischen Orte in der Bahn zum aufsteigenden Knoten, die Lage dieses Knotens zum Äquinoktialpunkt und die Neigung der Bahn gegen die Ekliptik (oder den Äquator) zu bestimmen. All dies lässt sich durch Auflösen eines einzigen sphärischen Dreiecks erreichen. Es seien $\AscendingNode$ die Länge des aufsteigenden Knotens, $i$ die Neigung der Bahn, $g$ und $g''$ die Breitenargumente bei der ersten und dritten Beobachtung und schließlich $l-\AscendingNode=h$, $l''-\AscendingNode=h''$. Bezeichnet nun in der vierten Abbildung $\AscendingNode$ den aufsteigenden Knoten, so sind die Seiten des Dreiecks $\AscendingNode AC$ gleich $AD'-\zeta,g,h$, und die jeweils gegenüberliegenden Winkel sind $i,180^{0}-\gamma,u$.
Wir erhalten also
\SourceLeaf{200}
\[
\begin{aligned}
\sin\tfrac12i\sin\tfrac12(g+h)&=\sin\tfrac12(AD'-\zeta)\sin\tfrac12(\gamma+u)\\
\sin\tfrac12i\cos\tfrac12(g+h)&=\cos\tfrac12(AD'-\zeta)\sin\tfrac12(\gamma-u)\\
\cos\tfrac12i\sin\tfrac12(g-h)&=\sin\tfrac12(AD'-\zeta)\cos\tfrac12(\gamma+u)\\
\cos\tfrac12i\cos\tfrac12(g-h)&=\cos\tfrac12(AD'-\zeta)\cos\tfrac12(\gamma-u).
\end{aligned}
\]

Die ersten beiden Gleichungen liefern $\tfrac12(g+h)$ und $\sin\tfrac12i$, die beiden übrigen $\tfrac12(g-h)$ und $\cos\tfrac12i$. Aus $g$ ergibt sich die Lage des Perihels zum aufsteigenden Knoten, aus $h$ die Lage des Knotens auf der Ekliptik. Schließlich ergibt sich $i$, wobei Sinus und Kosinus einander bestätigen. Zum selben Ergebnis gelangen wir mit dem Dreieck $\AscendingNode A''C''$, indem wir in den vorangehenden Formeln die Zeichen $g,h,A,\zeta,\gamma,u$ durch $g'',h'',A'',\zeta'',\gamma'',u''$ ersetzen. Zur weiteren Kontrolle der gesamten Arbeit ist es sinnvoll, die Rechnung auf beide Arten auszuführen. Ergeben sich dabei geringfügige Unterschiede zwischen den Werten von $i,\AscendingNode$ und der Länge des Perihels in der Bahn, so empfiehlt es sich, die Mittelwerte zu nehmen. Diese Unterschiede werden jedoch selten $0''{,}1$ oder $0''{,}2$ erreichen, sofern alle Rechnungen sorgfältig mit sieben Dezimalstellen durchgeführt wurden.

Wird schließlich anstelle der Ekliptik der Äquator als Grundebene gewählt, so ändert sich an der Rechnung lediglich, dass anstelle der Punkte $A,A''$ die Schnittpunkte des Äquators mit den Großkreisen $AB,A''B''$ zu nehmen sind.
